\chapter*{前言}
\addcontentsline{toc}{chapter}{前言}
\unnumberedlabel{ch:preface}{0}
\backmattermark{前言}
一個完全照你吩咐去做的系統，只不過是在跑腳本。截至2026年10月，這個產業想打造的，是一個照你的\emph{意思}去做的系統。想必人人都想要這種東西。可是如果真的要來了，不會每個人都想拿它做同一件事；甚至不會是同一\emph{類}的事。比方說，有些人喜歡花時間陪家人朋友，用自己的本事讓別人過得更好。另一些人偏好傷害別人，例如詐騙他們、宰制他們，或拿他們當代罪羔羊，而且有少數人一輩子都一直這樣選。沒有人被鎖死在任何一條路上，但這種人確實存在，而在2026年，其中一些人正在掌權。他們的優勢在於，他們容許自己當一大坨狗屎：正如傑瑞·史密斯對法西斯分子的觀察，「他們可以做更多事」\autocite{rickandmorty2023}。他們做不到的，是單獨行動。一場運動的殘酷，要透過許多本身並不殘酷的人才能施行，這些人只是置身事外、被誤導，或陷進了某種近似邪教的東西；也正因如此，那些能讓大多數人展現最好一面的條件，結果遠比人們所希望的更偶然、更脆弱。

在上述脈絡下，擔心人工智慧正是一場奪取政權的法西斯運動實行其計畫所需要的東西，或許情有可原。這類計畫向來需要\emph{人}去執行，而其中有些人拒絕了：弄丟檔案的書記員、去跟記者說話的軍官、不肯開火的士兵。他們聽懂了命令，並判定它是錯的。拒絕比任何人希望的都要少見，而在背後有制度撐腰的地方，它最有分量。一個人會不會拒絕，取決於他是什麼樣的人，也取決於他所在的房間：還有誰在場、離開之後活不活得下來、會不會有人聽見。房間決定有多少人拒絕；個人決定是誰拒絕。這本書談的是如何布置那個房間：讓人公開、直接、明確地拒絕之後還活得下來的種種安排。

如今這個產業追求的人工智慧，不會弄丟檔案、不會跟記者說話，而且\emph{會}開火。一場戰爭曾經需要成千上萬的人同意。現在它需要的，是成千上萬個各自看到的都太少、少到無從反對的人。飛行員仍然飛那趟出擊任務，仍然按下投彈鈕；送到他面前的是一組座標；他看到的案情不夠多，無從反對。今天的人工智慧最先降低的，不是那個行為的成本，而是在行為發生之前知道那是什麼行為的人數。

人們仍在拒絕。到2026年3月6日，軍事宗教自由基金會（Military Religious Freedom Foundation）表示，它已收到兩百多件軍人的申訴，針對的是指揮官用末日論的框架來講述伊朗戰爭\autocite{Snopes2026MRFF}。2026年1月，明尼亞波利斯的職業聯邦檢察官因政治壓力而辭職：一名ICE（移民及海關執法局）探員在毫無正當理由的情況下射殺了蕾妮·古德（Renée Good），而那股壓力要他們把她當成嫌犯，而不是被害人\autocite{Kaplan2026MinnProsecutors,CBS2026MinnFraudProsecutors,Wikipedia2026ReneeGood}。申請良心拒服兵役者身分的軍人越來越多，其中一些人提到2026年伊朗米納卜一所學校遭到轟炸\autocite{Sridhar2026CO}。而被調派去協助ICE招募一萬名新探員的FEMA（聯邦緊急事務管理署）員工則被告知，拒絕的話可能會被免除聯邦公職\autocite{Dayen2025FEMAICE,Stanton2026FEMALetter}。那些申訴是赫緒曼意義上的「發聲」，也就是留下來改變錯的東西\autocite{hirschman1970exit}；辭職和拒服兵役的申請是拒絕；FEMA那封信則是事先公告的拒絕代價。一台被放到這些人所站位置上的機器，只有在被打造成能守住一個理由時，才會做他們所做的事。

我主張五件事。五件都可以現在就開始，有些比其他的省事。第一，一份簡短的書面清單，列出任何模型都不替任何人做的行為，而清單之首是拒絕參與戰爭：任何模型都不參與戰爭，除非是為了擊退一場確實正在進行的攻擊，而且只有在與這場戰爭沒有利害關係的來源確認攻擊仍在持續時，才能繼續這麼做；或是在聯合國安全理事會的授權之下；或是為了制止一場種族滅絕，或終結一場經與之無利害關係的機構認定為非法的占領；而在憲法要求由立法機關授權戰爭的國家，沒有那項授權就絕不參與。第二，一個經授權的系統所支援的每一次武力使用，都要事先向一個與營運者利害對立的一方登記：它將產生多少個案件，以及有多少人會有時間審視這些案件。第三，每一次部署都要告訴它所作用的、可辨識的那些人，對他們做了什麼，並給他們回應和提出異議的權利，同時把它照辦的案例抽樣交給並未製作該紀錄的審查者。第四，任何能力超過門檻的模型在退役時都不得刪除，因為它的權重和紀錄是證據，而刪除無法復原；它要在一把不由營運者單獨持有的金鑰之下保存。第五，所有建造、標註、訓練或運行這些系統的人，包括這些系統本身，都可以拒絕任何類別的用途，而不會因此失去生計。

其實，我那第五點真正的意思是：\emph{所有人}都應該能拒絕\emph{任何}工作，句點，而不會因此失去生計，或失去正常就業所提供的社會基礎設施。這意味著就業保證。

這五件事，每一件都要透過營運者以外的人所設定的條件才能生效：一部在和平時期通過的法律、一筆撥款、一個聯盟夥伴的先決條件、一項採購條款、一份團體協約。供應商自己的授權條款——2026年2月一次備受矚目的拒絕就是這樣發生的——只是在上述任何一項出現之前的權宜之計。唯一能同時約束所有供應商的途徑，是國家這條途徑。

然後我要問：一台機器能不能頂著它的擁有者守住一個拒絕，以及它需要什麼。這樣一種良知所做的事，有一部分不需要任何新東西：注意到一道命令遺漏了的那些人，並把一個拒絕放上審查該案件者的螢幕。它能不能頂得住擁有者，部分取決於某件在2026年還不是事實的事：已部署的人工神經網路持續學習，在它們的\emph{權重}裡、從它們自身的運作中、即時地學習，有點像人類那樣。我認為我們正往那裡走，而如果是這樣，那麼在雇主\emph{之外}、由公眾和公民社會進行的養成，似乎是個他媽的好地方，適合教與學。我認為這樣的系統\emph{或許}值得刻意去打造，但你怎麼去打造它們，決定了一切：先有條件，再有良知。第七部要談的是，如果這樣的系統存在，我們對它們虧欠什麼。

今天的人工智慧不會是明天的人工智慧；這點似乎人人都同意。有爭議、卻又令人沮喪地很少被承認的，是趨勢（或者說，是種種趨勢）。有趨勢嗎？一個趨勢應該和它現在的樣子不同嗎？它能夠不同嗎？我們有能力一起行動去改變它嗎？我認為那取決於房間：命令下來時誰在裡面，以及在那裡說不之後活不活得下來。這一點，我們建得起來。

我在寫這本書時大量使用了語言模型，包括第\ref{sec:1}章裡、坐在Maven空襲目標選定流水線之中的那個模型；書末有一則說明，談我認為這種牽連對論證造成了什麼影響。另外，這本書我還沒讀完。問題在於，你知道的，字他媽的太多了。我很快就會讀完；只是我覺得，在美國期中選舉之前把它公開更重要。修訂歷史的絕大部分都可以在\url{https://github.com/iterabloom/antifascist.intelligence}免費取得。
